\documentclass[../../../include/open-logic-section]{subfiles}

\begin{document}

\olfileid{sth}{choice}{justifications}
\olsection{关于选择公理的内在考量}

那么，更广泛的问题是：良序原理、选择公理，甚至所有集合在大小上的可比性——不管是哪一个——能否得到辩护。

下面是一种\emph{内在}辩护的尝试。在 \olref[z][story]{sec} 中，我们引入了关于层级的若干原则。其中一条值得重申：

\begin{enumerate}
	\item[] \stagesacc{}。对于任意阶段 $S$，以及任意在阶段 $S$ \emph{之前}形成的集合：在阶段 $S$ 会形成一个集合，其成员恰好就是那些集合。在阶段~$S$ 不再形成其他任何东西。
\end{enumerate}

事实上，许多作者都曾提出，选择公理可以通过（类似于）这条原则来辩护。我们将简要说明这一思路。

我们将从一个简单的结论开始，它提供了选择公理的\emph{又一个}等价形式：

\begin{thm}[在 $\ZF$ 中]\ollabel{choiceset}
选择公理等价于下述原则：如果 $A$ 的元素互不相交且非空，则存在某个 $C$，使得对每个 $x \in A$，$C \cap x$ 都是单元素集。（我们称这样的 $C$ 为 $A$ 的{选择集}。）
\end{thm}

该结论的证明是直截了当的，留作习题。

\begin{prob}
证明 \olref[sth][choice][justifications]{choiceset}。如有困难，可在 \cite[pp.~242--3]{Potter2004} 中找到证明。
\end{prob}

关键在于，$A$ 的一个选择集恰好就是 $A$ 的某个选择函数的值域。因此，要为选择公理辩护，我们只需为其等价形式（即选择集的存在性）进行辩护即可。现在我们就来尝试这样做。

设 $A$ 的元素互不相交且非空。根据 \stageshier{}（参见 \olref[z][story]{sec}），$A$ 在某个阶段~$S$ 形成。注意，$\bigcup A$ 的所有元素在阶段 $S$ 之前就已经可用。现在，根据 \stagesacc{}，对于在~$S$ 之前形成的\emph{任何}集合，都会形成一个恰好以这些集合为成员的集合。换言之：每一个由先前可用集合构成的\emph{可能的}汇集，在~$S$ 都将存在。但是，选出能够构成~$A$ 的选择集的对象，这当然是\emph{可能的}；它只不过是 $\bigcup A$ 的某个非常特定的子集。因此：这样的选择集是存在的，满足要求。

以上便是基于内在理由为选择公理辩护的一个\emph{非常}简短的尝试。若要进一步探究这一想法，你应该阅读 Potter (\citeyear[\S14.8]{Potter2004}) 对此的精巧推演。

\end{document}